\chapter{Theory and Frameworks}


\section*{Notation}
\addcontentsline{toc}{section}{Notation}

\begingroup
\small
\renewcommand{\arraystretch}{0.9}
\setlength{\parskip}{0pt}
\begin{longtable}{@{}p{2.4cm}p{11.2cm}@{}}
\toprule
\textbf{Symbol} & \textbf{Meaning} \\
\midrule
\endfirsthead
\toprule
\textbf{Symbol} & \textbf{Meaning} \\
\midrule
\endhead
$X$ & Expression count matrix produced by the preprocessing pipeline, with features (transcripts or genes, depending on the pipeline output level) in rows and cells in columns \\
$x_{fc}$ & Observed integer count of feature $f$ in cell $c$ \\
$f$, $c$ & Feature and cell indices \\
$N$ & Number of cells \\
$R$ & Total number of sequencing reads \\
$L$ & Sequenced read length \\
$M$ & Mean number of reads per cell \\
$k$ & Length of the $k$-mers used to index the reference \\
$|\Sigma|$ & Alphabet size ($4$ for DNA) \\
$T$ & Number of annotated transcripts in the reference \\
$T_{obs}$ & Distinct transcripts detected in a single cell \\
$L_{avg}$ & Mean transcript length \\
$\theta_t$ & Relative abundance of transcript $t$ \\
$\ell_t$ & Effective length of transcript $t$ \\
$C_r$ & Equivalence class of transcripts compatible with read $r$ \\
$\hat{z}_{rt}$ & Fractional assignment of read $r$ to transcript $t$ \\
\midrule
\multicolumn{2}{@{}l}{\textit{Statistics}} \\
\midrule
$a_c$, $b_c$ & Binary presence indicators for a pair of features in cell $c$ \\
$O_{ij}$, $E_{ij}$ & Observed and expected counts in a $2 \times 2$ contingency table \\
$\mathrm{Row}_i$, $\mathrm{Col}_j$ & Marginal totals of that table \\
$\chi^2$ & Pearson chi-squared statistic \\
$C_{f,f'}$ & Co-expression coefficient between features $f$ and $f'$ \\
$p_{f,f'}$ & Independence $p$-value for that pair \\
\midrule
\multicolumn{2}{@{}l}{\textit{Detection}} \\
\midrule
$D_{f,S}$ & Differential expression (DEA) coefficient of feature $f$ in cluster $S$ \\
$S$ & Cell cluster \\
$\tau$ & Contrast threshold on the DEA difference within a gene \\
$\mathrm{gdi}_i$ & Cluster-specificity index of feature $i$ \\
\bottomrule
\end{longtable}
\endgroup

\section{Preprocessing Pipeline}

A preprocessing pipeline carries a single \gls{scRnaSeq} experiment from its data-generating mechanism to the \gls{countMatrix} consumed by every later analysis. That transformation is the bridge between raw sequencing output and statistical input, and it is organised as three successive steps: ingestion and quality control, mapping, and deduplication into a matrix. Ingestion and quality control turn raw reads into clean units, discarding the failures of the chemistry---empty droplets, multiplets, and low-quality molecules---so that the cells entering the analysis are genuine and distinct. Mapping assigns each read to the \gls{transcript} it originates from by comparing it against a reference transcriptome and reporting the set of transcripts with which it is compatible. Deduplication collapses the \gls{pcr} copies that descend from one captured molecule, grouping reads by cell barcode, UMI, and assigned transcript, so the count matrix records the number of distinct molecules per cell and feature rather than the number of reads.

Furthermore, the protocol chosen at the outset is not an incidental detail: it sets the depth, resolution, and sparsity of the whole experiment, and a different chemistry would reshape every phase and the matrix it ultimately produces. The subsections that follow take up each step in order, following a single experiment as it moves from raw sequencing output to the count matrix that every later analysis consumes. Figure~\ref{fig:pipeline} summarises the four phases and records what each one costs.

\begin{figure}[H]
\centering
\begin{tikzpicture}[
  node distance=7mm,
  box/.style={draw, rounded corners=1pt, align=center, text width=3.7cm, minimum height=0.8cm, inner sep=3pt, font=\small},
  note/.style={align=left, text width=6.2cm, font=\footnotesize, text=black!75},
  ar/.style={-{Latex[length=2mm]}, thick}
]
\node[box] (p0) {\textbf{0}\quad Protocol choice};
\node[box, below=of p0] (p1) {\textbf{1}\quad Ingestion and quality control};
\node[box, below=of p1] (p2) {\textbf{2}\quad Mapping};
\node[box, below=of p2] (p3) {\textbf{2.1}\quad Probabilistic resolution (EM)};
\node[box, below=of p3] (p4) {\textbf{3}\quad Deduplication and matrix assembly};
\node[note, right=8mm of p0] (n0) {Capture chemistry fixes which part of a transcript is ever read, and how many cells are affordable.};
\node[note, right=8mm of p1] (n1) {Empty droplets, multiplets and low-quality reads are discarded.};
\node[note, right=8mm of p2] (n2) {Each read is assigned to the transcript it originates from, against a reference transcriptome.};
\node[note, right=8mm of p3] (n3) {Pseudo-alignment only: transcripts sharing a terminal region become statistically indistinguishable.};
\node[note, right=8mm of p4] (n4) {What survives: one integer per feature and cell, the matrix $X$.};
\draw[ar] (p0) -- (p1);
\draw[ar] (p1) -- (p2);
\draw[ar] (p2) -- (p3);
\draw[ar] (p3) -- (p4);
\end{tikzpicture}
\caption{The four phases of the preprocessing pipeline and the information each one removes. Phase~2.1 exists only on the pseudo-alignment route, where ambiguous reads must be distributed probabilistically; direct alignment resolves every read by its coordinates and skips it. The chemistry chosen in Phase~0 bounds every later stage, so the losses accumulate rather than cancel.}\label{fig:pipeline}
\end{figure}

\subsection{Phase 0: Choosing the Protocol}

The design of any \gls{scRnaSeq} analysis pipeline is inherently constrained by the physical protocol used for data generation. This experimental choice dictates the scale, resolution, and \gls{sparsity} of the resulting \gls{countMatrix}. Table~\ref{tab:protocols} compares the droplet protocol against plate-based methods.

\begin{table}[h]
\centering
\small
\setlength{\belowcaptionskip}{10pt}
\caption{The two dominant single-cell library chemistries. Per-cell depth and cost trade off against throughput; total reads per day are set by the sequencer, not the chemistry. All figures describe a single sequencing run. Values are order-of-magnitude, after Ding et al.\ \cite{ding2020systematic}.}\label{tab:protocols}
\begin{tabular}{lcc}
\toprule
 & \textbf{Plate (SMART-seq2)} & \textbf{Droplet (10x)} \\
\midrule
Library preparation & $\sim 2$ days for $96$ cells & $\sim 1$ day for $10^{4}$--$10^{5}$ cells \\
Cells per experiment & $10^{2}$--$10^{3}$ & $10^{4}$--$10^{5}$ \\
Reads per cell & $\sim 10^{6}$ & $\sim 10^{4}$ \\
Detected transcripts per cell & $\sim 10^{4}$--$10^{5}$ & $\sim 10^{3}$ \\
Reads per day & $\sim 10^{8}$--$10^{9}$ & $\sim 10^{8}$--$10^{9}$ \\
Coverage & full transcript & 3' terminus only \\
Reagent cost per cell & $\approx 10$--$100\times$ higher & baseline \\
\bottomrule
\end{tabular}
\end{table}

The widespread adoption of droplet protocols is ultimately driven by economics. While a plate-based experiment costs one to five dollars per cell, droplet encapsulation profiles cells for a few cents, allowing researchers to assay two to three orders of magnitude more cells on a fixed budget \cite{conte2024opportunities}. This trade-off between per-cell depth and cell quantity is highly advantageous for standard tasks like cell-type clustering, marker discovery, and differential gene expression. For these gene-level questions, the statistical power gained by averaging over thousands of cells overcomes the severe sparsity within individual cells. However, the loss of transcript coverage remains a fundamental limitation for \gls{isoform}-level analyses.

\subsection{Phase 1: Data Ingestion and Quality Control}

The first computational challenge in a single-cell pipeline is purely one of scale and structure. Sequencers output raw data in \gls{fastq} format, representing each read across four distinct lines: a coordinate identifier\defn{Read identifier}{A line beginning with \texttt{@} that names the read and records the instrument, flow cell and coordinates it came from.}, the nucleotide sequence\defn{Nucleotide sequence}{The read itself, a string of A, C, G and T bases in the order the sequencer read them.}, a separator\defn{Separator}{A line beginning with \texttt{+} that marks the end of the sequence and may repeat the identifier.}, and the corresponding Phred-scale quality scores\defn{Per-base quality score}{One character per sequenced base, encoding the probability that the base was called incorrectly on the Phred scale.}. For droplet-based chemistries such as the 10x Genomics platform, this raw sequence carries critical multiplexing metadata that must be carefully isolated before any biological analysis begins. Specifically, the sequencer must extract a \gls{cellBarcode} that maps the read back to its discrete cell of origin, and a \gls{umi}---a random tag attached prior to \gls{pcr} amplification used to track individual captured mRNA molecules.

Before the reads can be quantified, they must be cleaned. Tools such as cutadapt\defn{cutadapt}{A read-trimming tool that removes adapter sequences and low-quality bases.} or fastp\defn{fastp}{A fast all-in-one read preprocessor that performs adapter trimming, quality filtering and reporting in a single pass.} trim synthetic adapters, discard low-quality bases, and extract the barcode and UMI payloads, while heuristic filtering separates valid cell barcodes from empty droplets and multi-cell collisions (multiplets). The step is optional, and is skipped under two conditions. When a vendor pipeline such as Cell~Ranger has already trimmed adapters and corrected barcodes and UMIs, repeating the pass costs a full read of the FASTQ file without altering a base. When the mapper is read-structure-aware and UMI-aware, as alevin-fry and Salmon are, residual adapters fall outside the indexed $k$-mers and are never matched, so correction is left to the quantifier. Skipping is safe only under those guarantees: on untrimmed reads handled by a mapper that does not correct UMIs, adapters mis-assign reads and sequencing errors inflate the number of distinct UMIs. For plate-based chemistry no barcode or UMI exists to extract and adapter content is low, so trimming reduces to a quality filter.

Because quality control touches every single nucleotide exactly once, it acts as a highly parallelisable streaming operation strictly bound by the total read count, $R$, and the read length, $L$. The computational time complexity is entirely linear. The physical footprint, by contrast, is large. One uncompressed FASTQ record---a $100$\,bp read together with its identifier and quality line---occupies roughly $250$ bytes, so a droplet experiment of $R \sim 10^{9}$ reads, the order of magnitude considered here, holds $\sim 2.5 \times 10^{11}$ bytes, about $250$\,GB before compression. Modern multi-core trimmers sustain throughputs on the order of $10^{8}$ reads per minute, so this I/O-bound filtering resolves in a few hours on a single workstation.

\subsection{Phase 2: Mapping Algorithms (Alignment vs.\ Pseudo-Alignment)}\label{sec:pseudo-align}

Quantifying transcript expression requires mapping short sequencing reads back to their reference transcripts of origin. The computational workload is primarily driven by the total read count, $R = N \cdot M$. For the two datasets analysed in this thesis, $R$ is approximately $10^{8}$ to $10^{9}$, establishing $10^{9}$ reads as a baseline scale (though pooled multi-sample atlases can easily reach $10^{11}$ reads).

The mapping process is further defined by three variables: the number of annotated isoforms ($T \sim 10^5$), the average transcript length ($L_{avg} \approx 1500\text{ bp}$), and the sequenced read length ($L \approx 90\text{--}150\text{ bp}$). In a $3'$ droplet assay, Read~1 isolates the cell barcode and UMI, leaving only Read~2 to capture the actual transcript sequence ($L$).

Traditional alignment algorithms\defn{Read alignment}{Base-by-base mapping of a read to the reference, giving exact coordinates at high computational cost.} (e.g., BWA \cite{li2009bwa}, Bowtie2 \cite{langmead2012bowtie2}) map these sequences base-by-base to assign exact genomic coordinates, typically relying on a Burrows-Wheeler Transform\defn{Burrows-Wheeler Transform (BWT)}{A reversible text transform that allows fast substring search in a compressed reference index.} index. While building this index is a one-time operation scaling as $\mathcal{O}(T \cdot L_{avg})$, querying it is computationally expensive. Because the algorithm must accommodate sequencing errors and mutations, it uses a branch-and-bound search that scales logarithmically with the reference size. Consequently, the total query time becomes:
$$ \mathcal{O}(T \cdot L_{avg}) + \mathcal{O}(R \cdot L \log T) $$
For the billions of reads generated in a single-cell experiment, this error-tolerant sequence searching becomes a severe bottleneck, making exhaustive alignment prohibitively slow.

Modern transcriptomics tools like Salmon, alevin-fry, and kallisto \cite{patro2017salmon, bray2016near, he2022alevin} bypass this bottleneck through pseudo-alignment. This approach relies on a key insight: downstream single-cell quantification only requires identifying the source transcript, not its exact genomic coordinates. By discarding base-level alignment, pseudo-alignment replaces logarithmic sequence searches with constant-time hash table lookups.

This approach operates in two computational phases. First, index construction decomposes the reference transcriptome into a colored de Bruijn graph\defn{Colored de Bruijn graph}{A graph of overlapping $k$-mers in which each $k$-mer is coloured by the transcripts containing it, used as the pseudo-alignment index.} that maps short sequences of length $k$ ($k$-mers) to their parent transcripts. Rather than materialising every theoretically possible sequence, the index only stores the $k$-mers actually observed in the reference. This maintains a manageable spatial footprint of $\mathcal{O}(d)$ for $d$ distinct $k$-mers; for a mammalian transcriptome $d \sim 10^8$, far below the $4^{31}$ theoretically possible $k$-mers. Second, during the query phase, each read is split into overlapping $k$-mers (roughly $60$ to $120$ per read). The algorithm queries the hash table for each $k$-mer and intersects the results to rapidly identify an equivalence class\defn{Equivalence class}{The set of transcripts that could have produced a read, since all contain the read's $k$-mers.}---the set of reference transcripts compatible with that read.

Because this method shifts the core operation from sequence alignment to hash retrieval, the per-read query cost decouples entirely from the size of the reference transcriptome. The total execution time simplifies to:
$$ \mathcal{O}(T \cdot L_{avg}) + \mathcal{O}(R \cdot L) $$

The practical impact of this reduction is substantial. Using the baseline scale of $R \sim 10^{9}$ reads, $L \sim 10^{2}\text{ bp}$, and $T \sim 10^5$ (where $\log_2 T \approx 16.6$), traditional alignment's query phase requires approximately $1.7 \times 10^{12}$ operations. In contrast, pseudo-alignment requires roughly $1.0 \times 10^{11}$ operations. This $50$- to $100$-fold performance leap occurs because pseudo-alignment drops the asymptotic $\log_2 T$ factor and replaces expensive string matching with efficient hash lookups, ultimately reducing processing times from days to hours.

\subsubsection{Phase 2.1: Probabilistic Resolution via Expectation-Maximization}

Pseudo-alignment introduces a distinct computational challenge: multi-mapping reads\defn{Multi-mapping read}{A read compatible with more than one transcript because those transcripts share its sequence.}. These occur when a read's $k$-mers match multiple transcripts that share exonic sequences. If a read falls entirely within a shared exon, it carries no sequence information to identify its specific isoform of origin, prompting the pseudo-aligner to return a set of valid candidates rather than a single definitive transcript. This ambiguity is particularly severe in $3'$-biased protocols, where short reads exclusively capture the terminal region and frequently match every isoform sharing that terminal exon \cite{galfre2021cotan}.

Discarding these ambiguous reads is statistically unsound. Because they constitute the vast majority of alignments in $3'$-biased data, removing them would severely deplete the count matrix and artificially bias isoform proportions toward rare variants that happen to carry unique sequences. Instead, modern pseudo-aligners resolve this ambiguity by framing quantification as maximum-likelihood estimation over a generative model of the sequencing experiment \cite{patro2017salmon}.

In this model, each read fragment is assumed to be produced in two stages. First, a transcript $t$ is drawn with probability proportional to its abundance $\theta_t$. Second, a start position is drawn uniformly from the $\ell_t$ valid positions where a fragment of the observed length could begin. This admissible count, $\ell_t$, represents the transcript's \defn{Effective length}{The number of fragment start positions a transcript can generate; approximately its length minus the fragment length, plus one. It corrects for the fact that longer transcripts yield more fragments at equal abundance.}. Crucially, the true transcript of origin remains latent; the pseudo-aligner only reports the equivalence class $C_r$ \cite{bray2016near}. Marginalising over the unobserved origin yields the probability of observing that class under a given abundance vector:
$$ P(C_r \mid \theta) = \frac{\sum_{t \in C_r} \theta_t \, \ell_t}{\sum_{t'} \theta_{t'} \, \ell_{t'}} $$

Because the origin of each read is unobserved, this likelihood is optimised using the Expectation-Maximization (EM)\defn{Expectation-Maximization (EM)}{An iterative algorithm that estimates transcript abundances and resolves ambiguous read assignments jointly until they converge.} algorithm, which alternates between two phases until convergence. In the Expectation step (E-step)\defn{E-step / M-step}{The two alternating EM phases: distribute reads by the current abundance estimate (E), then re-estimate abundances (M).}, the current abundance estimate is used to compute the posterior probability that read $r$ originated from transcript $t$, given that it belongs to the class $C_r$:
$$ \hat{z}_{rt} = \frac{\theta_t \, \ell_t}{\sum_{t' \in C_r} \theta_{t'} \, \ell_{t'}} \quad \text{for } t \in C_r, $$
and $\hat{z}_{rt} = 0$ otherwise. The effective-length term $\ell_t$ is critical here: it accounts for the fact that a longer transcript contributes proportionally more fragments at equal abundance, ensuring $\hat{\theta}$ is not biased toward shorter isoforms. The algorithm then proceeds to the Maximization step (M-step), where each abundance is re-estimated as the total posterior weight assigned to its transcript:
$$ \theta_t \propto \sum_{r \,:\, t \in C_r} \hat{z}_{rt} $$
These values are subsequently normalised so that $\theta_t$ sums to exactly one across the transcriptome.

Two properties of this iteration influence the interpretation of the resulting counts. First, because each EM round is guaranteed not to decrease the log-likelihood, the estimates converge monotonically to a local optimum, though not necessarily a global one \cite{dempster1977maximum}. Since the likelihood surface is frequently flat or multimodal, implementations typically initialise $\theta$ uniformly and stop when the relative change in log-likelihood falls below a specified tolerance. Second, the model relies on a closed-world assumption: every read is assumed to originate from an annotated transcript. Consequently, incomplete references or unannotated isoforms propagate directly into $\hat{\theta}$. The resulting point estimate lacks intrinsic uncertainty, meaning downstream methods that require variance---such as the differential transcript usage (DTU) tools discussed later---must recover it by bootstrapping the quantification step \cite{love2018swimming}.

Computationally, because the EM updates run over equivalence classes rather than individual reads, the algorithm scales with the number of distinct classes. Tools like alevin-fry's parsimonious variant exploit this by collapsing transcripts that share identical $k$-mer sets into single groups, effectively performing gene-level quantification where isoform resolution is impossible \cite{he2022alevin}.

This collapsing behaviour highlights a fundamental identifiability limit inherent to $3'$ capture protocols. Consider two isoforms, $t$ and $t'$, of the same gene whose reads fall entirely into equivalence classes containing both---a common scenario when a shared terminal region covers their effective lengths. During the E-step, posterior weight is assigned solely through the ratio $\theta_t \ell_t / (\theta_t \ell_t + \theta_{t'} \ell_{t'})$. If $\ell_t \approx \ell_{t'}$, this ratio reduces to $\theta_t / (\theta_t + \theta_{t'})$. As a result, the likelihood depends on the two abundances only through their sum and remains flat along $\theta_t + \theta_{t'} = \text{constant}$. Every proportional split of abundances explains the observed reads equally well. This identifiability problem is bounded by the sequence divergence at the terminal region, not by the estimator, the reference, or sequencing depth; sequencing the same fragment more deeply cannot separate two transcripts that the fragment itself cannot distinguish. Consequently, transcript-level analyses of $3'$-biased data can only reliably recover isoform switches that differ at or near their $3'$ ends.

\subsection{Phase 3: Deduplication and Matrix Generation}

Once each read has been assigned to a transcript, the remaining task is to count molecules rather than reads. The protocol attaches a \gls{umi}---a short random sequence---to each captured molecule before amplification, so every read carries both a \gls{cellBarcode} identifying its cell and a UMI identifying its molecule. Because PCR copies inherit both tags, the reads descending from one molecule are indistinguishable duplicates, and collapsing each such group recovers the molecule counts that sequencing had inflated.

The collapse must be performed on barcode, UMI \emph{and} assigned transcript together. A UMI holds only a few random bases, so within one cell two unrelated molecules can share a UMI by chance; grouping on barcode and UMI alone would merge them, whereas including the transcript keeps them apart whenever they originate from different features.

The result is the expression count matrix $X$, whose entry $x_{fc}$ records the number of distinct molecules observed for feature $f$ in cell $c$. Amplification no longer distorts the counts, so the differences remaining between cells reflect biology and capture efficiency rather than how often each molecule happened to be copied.



\section{COTAN Framework}\label{sec:cotan-framework}

The transcript-level matrix $X$ produced above acts as the natural input for the statistical framework presented here, yet it simultaneously introduces the primary analytical difficulty: the shallow per-cell depth heavily imposed by the droplet hardware leaves the vast majority of its entries at zero. The COTAN framework, introduced by Galfr\`e et al.\ \cite{galfre2021cotan}, was specifically engineered to analyse this zero-heavy structure directly, and has since been logically extended to recover entire cell clusters purely from co-expression patterns \cite{galfre2026gene}. The following subsections formalise this sparsity and detail the contingency-table machinery COTAN employs to successfully bypass traditional density-based statistics.

\subsection{The Sparsity Problem}\label{sec:the-sparsity-problem}

Single-cell count matrices exhibit extreme structural sparsity. The disparity between the reference transcriptome and the yield of a single cell sets the scale: if $T \sim 10^5$ annotated isoforms sit in the reference while a single cell recovers only $T_{obs} \sim 10^3$ distinct transcripts, raw alignment space carries an initial zero rate of approximately $99\%$. Standard preprocessing pipelines subsequently filter out unexpressed or ultra-rare features, condensing the effective feature space to $T_{filtered} \approx 1.2 \times 10^4$--$2.3 \times 10^4$ viable transcripts, yet the observed zero rate typically remains between $70\%$ and $90\%$ \cite{galfre2021cotan}.

These zeros fall into two categories: biological zeros\defn{Biological zero}{A zero count reflecting genuine absence of a transcript in a cell, as opposed to a technical dropout.} (where the cell did not express the transcript) and technical \glspl{dropout} (where the molecule was present but failed to be captured). Because the cumulative efficiency of droplet capture, reverse transcription, and sequencing is only $10\%$ to $20\%$ \cite{jiang2022statistics}, technical dropouts dominate. Once recorded, both types of zeros are numerically identical in the count matrix and cannot be distinguished by count-based tests alone: a transcript that was truly absent and one that was merely missed leave the same mark \cite{sarkar2021separating}. Furthermore, the $3'$ terminus\footnote{\textbf{3' (three-prime) terminus} is the end of a transcript at which the poly-A tail is added; capture chemistry anchored here reads only the terminal fragment, not the transcript body.} \gls{tagBias} exacerbates this sparsity by collapsing distinct isoforms that share terminal exons into a single observational unit. That ambiguity, rather than the sparsity itself, is the fundamental statistical barrier for conventional tools designed for deep, dense \gls{bulkRnaSeq} data.

Applying standard density-based statistical methods to such sparse matrices introduces severe, compounding computational artifacts. Log-normalisation, which strictly involves adding a pseudo-count before logarithmic scaling, heavily distorts low-abundance measurements and artificially compresses legitimate zeros into spurious scaling patterns \cite{galfre2021cotan}. Standard linear or rank-based association measures, such as Pearson or Spearman correlations, fail entirely in this domain, as they inevitably end up measuring the coincidence of technical dropouts rather than identifying true regulatory co-expression \cite{galfre2021cotan}. Furthermore, attempting to artificially densify the matrix through heuristic imputation---specifically replacing zeros with smoothed estimates derived from neighbouring cells---frequently introduces fabricated correlations and severely compromises sensitive downstream signals \cite{galfre2021cotan}. To fundamentally circumvent these pitfalls, COTAN explicitly abandons imputation and variance-stabilising transformations entirely, choosing instead to evaluate the zeros as a direct signal through strict binary discretisation.

\subsection{Contingency Table Solution}

As stated, COTAN cleanly bypasses the flawed continuous count amplitudes by evaluating pairwise feature dependencies using asymptotic independence testing applied strictly to binary presence and absence states \cite{galfre2021cotan}. For any given feature pair $(f, f')$ across $N$ cells, the continuous counts are logically reduced to discrete indicator variables for each cell $c$:
$$ a_c = \mathbf{1}[x_{fc} > 0], \qquad b_c = \mathbf{1}[x_{f'c} > 0] $$
where $a_c$ marks whether feature $f$ is detected in cell $c$ and $b_c$ whether its partner $f'$ is. Letting the discrete states be indexed by $i, j \in \{0,1\}$, the framework constructs a comprehensive $2 \times 2$ \gls{contingencyTable} where $O_{ij}$ directly denotes the observed number of cells exhibiting state $a_c = i$ and $b_c = j$. The marginal totals for this specific table are defined as $\mathrm{Row}_i = \sum_j O_{ij}$ and $\mathrm{Col}_j = \sum_i O_{ij}$, mathematically constrained by the total cell count $\sum_i \mathrm{Row}_i = \sum_j \mathrm{Col}_j = N$. Under the standard null hypothesis of statistical independence, the expected cell count $E_{ij}$ is calculated directly as:
$$ E_{ij} = \frac{\mathrm{Row}_i \cdot \mathrm{Col}_j}{N} $$

Table~\ref{tab:contingency} lays the four counts out explicitly. The two diagonal cells hold cells in which the features are detected together ($O_{11}$) or neither is detected ($O_{00}$); the off-diagonal cells hold cells defining the difference between them, in which exactly one of the two is present. Co-expression resides in the tendency of the count to fall on a diagonal rather than on the anti-diagonal, and the expected value $E_{ij}$ above is what that count would be under the independence hypothesis.

\begin{table}[htbp]
\centering
\small
\setlength{\tabcolsep}{7pt}
\caption{The $2 \times 2$ contingency table built for a feature pair $(f, f')$ across $N$ cells. Entry $O_{ij}$ counts cells with $a_c = i$ and $b_c = j$; the margins are the feature's detection totals.}\label{tab:contingency}
\begin{tabular}{lccc}
\toprule
 & \textbf{$f'$ detected ($b_c = 1$)} & \textbf{$f'$ absent ($b_c = 0$)} & \textbf{Row total} \\
\midrule
$f$ detected ($a_c = 1$) & $O_{11}$ & $O_{10}$ & $\mathrm{Row}_1$ \\
$f$ absent ($a_c = 0$) & $O_{01}$ & $O_{00}$ & $\mathrm{Row}_0$ \\
\midrule
Column total & $\mathrm{Col}_1$ & $\mathrm{Col}_0$ & $N$ \\
\bottomrule
\end{tabular}
\end{table}
Deviation from this independent expectation is evaluated using Pearson's chi-squared statistic \cite{galfre2021cotan}:
$$ \chi^2 = \sum_{i \in \{0,1\}} \sum_{j \in \{0,1\}} \frac{(O_{ij} - E_{ij})^2}{E_{ij}} $$

The result is reported not as a statistic but as a single coefficient. COTAN defines the co-expression coefficient $C_{f,f'}$ by rescaling the deviation so that it lies in $(-1, 1)$ and, under the null hypothesis, satisfies $N\,C_{f,f'}^2 \sim \chi^2(1)$ over the $N$ cells of the table. The pairwise $p$-value, the quantity used for screening throughout this thesis, then follows directly:
$$ p_{f,f'} = P\!\left(\chi^2(1) > N\,C_{f,f'}^2\right) $$
A positive coefficient means the two features are detected together more often than independence predicts; a negative one means less often, which is the mutual exclusivity that Chapter~\ref{chap:methods} treats as the signal of interest.

One modification separates $C_{f,f'}$ from the textbook Phi coefficient of a $2 \times 2$ table, and it lies in the expectation rather than in the statistic. Taking $E_{ij}$ proportionally to the marginals, as above, assumes that every cell has the same chance of expressing a given feature. Cells differ widely in library size, so the cells that capture less RNA accumulate zeros across all features at once, and a proportional expectation would read that shared under-detection as co-expression and inflate the coefficient. COTAN therefore derives the expected co-occurrence from a model of the zeros themselves, calibrated per feature and per cell from the zeros actually observed; the reader is referred to the COTAN paper for the exact form of that model \cite{galfre2021cotan}. Co-expression is thus measured against how often two features' zeros were expected to coincide, not against their raw marginals.

The chi-squared approximation behind $p_{f,f'}$ is asymptotic in the number of cells, not in the counts, and the distinction matters here. With $N$ in the thousands the calibration is reliable for features detected in a reasonable fraction of cells, which covers most of a standard gene matrix. It degrades in the opposite regime, where a feature is rare enough that several expected cell counts $E_{ij}$ fall below about five; Pearson's statistic is then anti-conservative, and such a pair is best read as inconclusive rather than significant. Filtering cannot remove this effect for the rarest transcripts.

\subsection{Output Metrics}

The framework reduces the machinery above to three output metrics \cite{galfre2021cotan}.

The co-expression matrix $C$ captures the pairwise associations between features through the calculated coefficient $C_{f,f'}$. Within this matrix, positive values represent co-expression, whereas negative values indicate mutual exclusivity.

The \gls{gdi} is calculated from the co-expression matrix prior to any cellular clustering and functions as a per-feature measure of cluster specificity. For feature $i$ it takes the mean of the smallest $5\%$ of the feature's pairwise $p$-values, $\bar{p}_i$, and applies the transformation $\mathrm{gdi}_i = \ln(-\ln \bar{p}_i)$ to compress the distribution into an interpretable scale \cite{galfre2021cotan}. Features with constitutive expression form a uniform lower tail around a null median of roughly $1.305$. Conversely, features that vary across subpopulations produce very small $p$-values and rise above approximately $1.5$, which separates cell-type markers from housekeeping features.

Differential expression analysis (DEA) determines whether a single feature $f$ is statistically enriched within a specific cell cluster $S$. It applies the contingency-table structure to compare the presence or absence of feature $f$ against cells located inside versus outside cluster $S$, yielding the correlation coefficient $D_{f,S}$ on the bounded scale $(-1, 1)$. A positive $D_{f,S}$ reveals that the feature is detected more frequently inside the cluster, while a negative value signifies it is detected less often. Because all DEA coefficients share this standardized scale, the enrichment values can be directly compared across different features and clusters.

COTAN also provides functions for cell clustering, separate from these metrics but built on the GDI: a uniform-clustering step that groups cells into clusters whose gene expression is homogeneous, checked by ensuring no more than a small fraction of genes exceed a GDI threshold, followed by a merging step that collapses clusters which remain uniform when joined. DEA operates on whichever clusterization is supplied, so the two are connected but distinct capabilities.

\section{The Theoretical Basis of DTU in Sparse scRNA-seq}

\subsection{The Probabilistic Definition of a Switch}

Differential Gene Expression (\gls{deg}) and Differential Transcript Usage (\gls{dtu}) target fundamentally different mathematical properties of the transcriptome. DGE strictly measures overall changes in total, aggregated gene abundance between conditions; if $\theta_g$ represents the total absolute abundance of a parent gene $g$, DGE models shifts in $\theta_g$ directly. DTU, conversely, models the conditional probability of observing a specific transcript $t$ given that its parent gene is expressed. Let $\theta_t$ represent the absolute abundance of transcript $t$; the continuous relative proportion of that transcript is therefore $P(t \mid g) = \theta_t / \theta_g$. A theoretical isoform switch occurs when the conditional probability $P(t_1 \mid g)$ of one variant significantly increases between biological conditions or cellular states while $P(t_2 \mid g)$ simultaneously decreases, a regulatory shift that can occur entirely independently of any change in the total aggregated abundance $\theta_g$.

\subsection{Translating Continuous Proportions to Binary Sparsity}

Traditional bulk RNA-seq tools model the continuous proportion $P(t \mid g)$ directly, a statistical approach that relies on deep, junction-spanning coverage to provide stable density. In the extreme sparsity of droplet data (Section~\ref{sec:the-sparsity-problem}), a continuous proportional model breaks down. Because the cumulative efficiency of droplet capture and sequencing is exceedingly low, as an isoform's relative proportion drops, its already-low integer count is overwhelmingly pushed entirely into the zero state, registering as \gls{dropout}. Consequently, a continuous regulatory shift in biological proportions manifests mathematically in the single-cell count matrix as asymmetric binary presence and absence across cells.

\subsection{The Theoretical Limit of Observability}

The underlying pseudo-alignment quantification model imposes a strict theoretical boundary on which DTU events are mathematically identifiable from the matrix. As established in Section~\ref{sec:pseudo-align}, the pseudo-aligner maps sequences not to exact coordinates but to equivalence classes $C_r$, representing the set of all transcripts compatible with a given read. If two actively switching isoforms $t_1$ and $t_2$ differ strictly in their internal exonic structure but share identical sequences at their $3'$ terminus, all terminal reads generated by droplet protocols will invariably fall into an equivalence class containing both variants. During quantification, the Expectation-Maximization (EM) algorithm assigns posterior weights based on the length-adjusted ratio $\theta_{t_1} \ell_{t_1} / (\theta_{t_1} \ell_{t_1} + \theta_{t_2} \ell_{t_2})$. If their effective lengths are approximately equal ($\ell_{t_1} \approx \ell_{t_2}$), the likelihood surface remains entirely flat along the constraint $\theta_{t_1} + \theta_{t_2} = \text{constant}$. Every proportional split of abundances explains the observed terminal reads equally well, rendering the true $P(t \mid g)$ unidentifiable. Therefore, the framework is theoretically constrained: it can only successfully detect DTU events---such as alternative $3'$ terminus usage or terminal exon variation---where sufficient sequence divergence allows the short terminal fragments to map to distinct, informative equivalence classes.

\subsection{Mutual Exclusivity Signal}\label{sec:mutual-exclusivity}

Within the adapted COTAN framework, negative statistical associations---formally defined as \gls{mutualExclusivity}---between features serve as the foundational indicator of alternative usage patterns. When applied specifically to multiple isoforms of the exact same parent gene, this structural mutual exclusivity is biologically consistent with alternative splicing or \gls{isoformSwitch} events, wherein different cellular states preferentially express entirely distinct transcript variants.

However, mutual exclusivity in isolation is insufficient to prove a functional switch, since a strong negative association can also arise from sparse detection rather than from a real switch. One way to demonstrate the switch, consistent with the $P(t \mid g)$ definition, is bidirectional opposition: a cluster in which one isoform is preferred and another in which the preference flips to its partner. This confirms that different cell types genuinely use different isoforms, rather than leaving the exclusivity explainable by sparsity.